% appx-hw-F1-summary

\section{Appendix F.1 --- Summary of Realisation}

\textit{One IP, three SKUs, one shuttle. The TRI-1 IP defined in the manuscript appendices and the PhD monograph is now physically taped out in three configurations.}

On \textbf{2026-05-19} (shuttle close 06:59 +07), the TRI-1 architecture was submitted in three configurations to the TinyTapeout \textbf{TTSKY26b} open-source SkyWater SKY130 shuttle. All three projects share the same root anchor identity $\varphi^{2}$ + $\varphi^{-2}$ = 3, the same Theorem 36.1 cross-die anchor \_\_C0\_\_, and the same Apache-2.0 licence umbrella.

\subsection{Three-SKU mapping (Phi / Euler / Gamma)}

\begin{tabular}{|l|l|l|l|l|}
\hline
SKU \& Top module \& Tile \& Project \# \& Anchor byte \\
\hline
Phi (Nano) \& \_\_C1\_\_ \& 1$\times$1 \& \#4914 \& \_\_C2\_\_ \\
Euler (Mid) \& \_\_C3\_\_ \& 8$\times$2 \& \#4915 \& \_\_C4\_\_ \\
Gamma (Max) \& \_\_C5\_\_ \& 8$\times$4 \& \#4913 \& \_\_C6\_\_ \\
\hline
\end{tabular}
\textit{Cross-die anchor on reset:} \_\_C7\_\_ (Theorem 36.1).

\textbf{Aggregate counts.} 246 RTL modules total across the three SKUs; all under Apache-2.0; no proprietary IP; no foundry NDA-bearing collateral.

\textbf{Process.} SkyWater 130 nm open PDK via TinyTapeout TTSKY26b shuttle (\_\_C8\_\_).

\textbf{Provenance DOI.} Trinity v1.0.0 framework deposited as \_\_C9\_\_ (cited in Sections 8.4, 12.1, and Appendix A.1).

\subsection{Honest performance window}

The following figures are \textit{projected at slow-slow corner on SKY130 with the TT shuttle harness}. They are not measured silicon; measurement awaits shuttle return.

- Throughput: \textbf{~1 GOPS} ternary at \textbf{~50 MHz} at \textbf{~1 W}.
- Energy figure of merit reported as a projected window only.
- No claim of competitive parity with commercial AI accelerators is made or implied.